% 本文件由 LaTeX 排版 Agent 根据有效 translation.json 自动生成。
% author=John Chrysostom; work_id=2309; title=费肋孟书讲道集

\chapter{费肋孟书讲道集}

% structure: p0001 source=h1 target=omitted reason=page-title-already-rendered-by-parent

论题\newline{}讲道一\newline{}讲道二\newline{}讲道三\par

\SourceRecord{Translated by Philip Schaff.；Nicene and Post-Nicene Fathers, First Series；Vol. 13.；Edited by Philip Schaff.；Buffalo, NY: Christian Literature Publishing Co.,；1889.；<http://www.newadvent.org/fathers/2309.htm>.}

% structure: p0001 source=h1 target=section reason=page-title

\section{要旨}

首先，必须陈明这封书信的要旨，然后再讨论其中受到质疑的问题。那么，要旨是什么呢？费肋孟是一位令人钦佩且品格高尚的人。他是一位令人敬佩的人，这从以下事实显而易见：他的全家都是信徒，而且这些信徒如此之多，以致被称为一个教会；因此他在此书信中说：\BibleQuote{以及在你家的教会。}\BibleRef{费 1:2} 他也见证了他的大服从，以及\BibleQuote{众圣人的心肠因你而舒畅。}\BibleRef{费 1:7} 并且他自己在此书信中吩咐他为他预备住处。\BibleRef{费 1:22} 因此在我看来，他的房子完全是圣人们的接待之所。这位优秀的人当时有一个名叫敖尼息摩的奴隶。这个敖尼息摩偷了主人的东西后逃跑了。因为他偷了东西，请听他所说的话：\BibleQuote{他若亏负了你，或欠你什么，都归在我账上。}\BibleRef{费 1:18-19} 因此他来到罗马的保禄那里，在狱中找到了他，享受了他的教导之恩，并在那里领受了圣洗。因为他在那里获得了圣洗的恩赐，这从他的话中显明：\BibleQuote{就是我在锁链中所生的。}\BibleRef{费 1:10} 因此保禄写信，将他推荐给他的主人，要求他无论如何要宽恕他，并接纳他如同一个重生的人。\par

但有些人说，将这封书信附录是多余的，因为他只是为一个人请求一件小事。让提出这些反对意见的人知道，他们自己应受许多谴责。因为不仅这些为如此必要之事所写的小书信应当被收录，而且我希望能够遇到一个人，他能向我们讲述宗徒们的历史，不仅包括他们所写和所说的一切，还包括他们其余的生活，甚至他们吃什么、何时吃、何时行走、坐在哪里、每天做什么、在哪些地方、进了哪座房子、住在哪里——详尽地叙述一切，因为他们所做的一切都充满了益处。但大多数人不知道由此产生的益处，反而加以指责。\par

因为如果我们仅仅看到他们坐过或被囚禁的地方，那些无生命的场所，就常常将我们的思想带到那里，想象他们的德行，被它激发，变得更加热心；那么，如果我们听到他们的话语和其他行为，就更会如此。但关于一位朋友，人会询问他住在哪里、做什么、去哪里；那么，我们难道不应该对世界这些普遍的导师也做这样的询问吗？因为当一个人过属灵生活时，这种人的习惯、走路、言语和行为，总之，与他有关的一切，都有益于听众，没有任何东西是阻碍或妨碍。\par

但这对你是有益的：要知道这封书信是就必要的事发出的。所以请看，藉此有多少事情得到纠正。首先，我们得到这一点：凡事都应当认真。因为如果保禄对一个逃跑的、偷窃的、强盗般的奴隶如此关心，不拒绝也不羞于用这样的推荐把他送回去；那么，我们更不应当在这样的事上疏忽。第二，我们不应该放弃奴隶这一族群，即使他们已陷入极端的邪恶。因为如果一个偷窃逃跑的人变得如此有德，以致保禄愿意他作自己的同伴，并在此书信中说：\BibleQuote{他可以在你位上服事我。}\BibleRef{费 1:13} 那么，我们更不应当放弃自由人。第三，我们不应该将奴隶从他们主人的服务中撤走。因为如果保禄对费肋孟如此有信心，却仍不愿未经主人同意而留下敖尼息摩——尽管他如此有用且有益于服事自己——那么，我们更不应当这样做。因为如果仆人如此优秀，他无论如何应当继续在那服事中，并承认主人的权柄，以便他能成为那家中众人的益处。你为什么从灯台上取下灯，放在斗底下呢？\par

我希望能将那些在外面的（仆人）带进城来。你说：\Quote{如果他也会败坏呢？} 请问，为什么他会败坏？因为他进了城？但请想想，在外面他会更加败坏。因为人在里面尚且败坏，在外面就会更加败坏。因为在这里，他会被免除必要的挂虑，主人替他承担这挂虑；但在那里，对这些事的挂虑或许会把他从更必要更属灵的事上拉开。为此，有福的保禄在给他们最好的劝告时说：\BibleQuote{你蒙召时是奴隶吗？不要为此挂虑；但即使你能得以自由，也宁可守住奴隶身份。}\BibleRef{格前 7:21} 也就是说，要停留在奴役中。但比这一切更重要的是，天主的话不被亵渎，正如他自己在一封书信中所说：\BibleQuote{凡负轭为奴的，应认为自己的主人堪受一切尊敬，以免天主的名和祂的教训被亵渎。}\BibleRef{弟前 6:1} 因为外邦人也会说，即使是一个奴隶也能中悦天主。但现在许多人被迫亵渎，并说基督教被引入生活是为了颠覆一切，主人被夺走他们的仆人，这是一件暴力的事。\par

我还要说另一件事。他教导我们，如果我们的家人有德，就不要以他们为耻。因为如果保禄，最令人钦佩的人，如此为这个奴隶说话，那么我们更应当为我们自己的仆人说话。既然有这么多益处——我们还没有列举全部——有人会认为这封书信被收录是多余的吗？这岂不是极端的愚蠢吗？那么，恳求你们，让我们专心致志于宗徒所写的这封书信。因为我们已从中获得如此多的益处，我们将从经文中获得更多。\par

\SourceRecord{Translated by Philip Schaff.；Nicene and Post-Nicene Fathers, First Series；Vol. 13.；Edited by Philip Schaff.；Buffalo, NY: Christian Literature Publishing Co.,；1889.；<http://www.newadvent.org/fathers/23090.htm>.}

% structure: p0001 source=h1 target=section reason=page-title

\section{论费肋孟书 第一篇讲道}

\BibleRef{费 1:1-3}\par

\BibleQuote{保禄，为耶稣基督被囚的，以及兄弟弟茂德，致我们亲爱的同工费肋孟，并致我们亲爱的亚腓亚，以及我们战友阿尔基颇，并致在你家里的教会：愿恩宠与平安由天主我们的父和主耶稣基督赐与你们。}\par

这些话是替一个仆人对主人说的。一开始，他就打击他的傲气，不容他感到羞耻；他平息了他的怒气。他自称囚犯，以此打动他，使他收敛，并使他看出现世之事毫无价值。因为若为基督的缘故带锁链不是羞耻而是夸耀，那么奴役更不应被视为耻辱。他说这话并非高抬自己，而是出于好意，为表示他值得信任；这样做也不是为自己，而是为了更容易获得恩准。仿佛他在说：\Quote{因着你的缘故，我戴上了这锁链。}正如他也在别处说过，那里显示了他的关怀，而这里则显示了他的可信赖。\par

没有什么比被称为\Quote{基督的烙印}更大的夸耀了。\BibleQuote{因为我身上带着耶稣的烙印。}\BibleRef{迦 6:17}\par

\Quote{主的囚犯。}因为他是为主而被捆绑的。谁听到基督的锁链而不敬畏，不自谦呢？谁不准备舍命，更何况一个家仆呢？\par

\BibleQuote{以及兄弟弟茂德。}\par

他又联合另一人与自己一起，使得那人因受多人恳求而更易让步并给予恩惠。\par

\BibleQuote{致我们亲爱的同工费肋孟。}\par

若\Quote{亲爱的}，则他的坦言并非鲁莽或冒昧，而是大爱的明证。若\Quote{同工}，则他不仅可在此事上受教导，更应将此视为恩惠。因为他在满足自己，他在建立相同的工作。所以，无需任何请求，他说，你有另一种必须给予恩惠的理由。因为如果他对福音有益，而你也渴望促进福音，那么你不应被恳求，反而应恳求别人。\par

\BibleRef{费 1:2} \BibleQuote{并致我们亲爱的亚腓亚。}\par

在我看来，她是他的生活伴侣。请注意保禄的谦卑：他不但请弟茂德与他一同求情，而且不仅请求丈夫，还请求妻子，以及另一位可能是朋友的人。\par

\BibleQuote{以及阿尔基颇，我们的战友。}\par

他不愿以命令的方式成就这些事，也不因对方未立即应允而生气；反而恳求他们做一件陌生人也可能为援助他的请求而做的事。因为不仅受多人请求，而且请求被向多人提出，都有助于请求获准。因此他说：\Quote{\BibleQuote{以及阿尔基颇，我们的战友。}}若你是战友，你也应在这些事上操心。这位阿尔基颇就是他在\WorkTitle{哥罗森书}中所说的：\Quote{\BibleQuote{你们要对阿尔基颇说：务要留心你在主里所受的职务，务要把它完成。}}\BibleRef{哥 4:17}依我看来，他请来一同请求的这个人也是圣职界的一员。他称他为战友，好使他一定与他合作。\par

\BibleQuote{并致在你家里的教会。}\par

在这里，他甚至连奴隶也没有忽略。因为他知道，常常连奴隶的话语也能推翻其主人；尤其当他的请求是为一个奴隶时。也许正是他们特别激怒了主人。因此他不使他们陷入嫉妒，而是藉着在问候中与主人并列来荣耀他们。他也没有让主人生气。因为若他提名提及他们，主人可能发怒。若他完全不提他们，主人也可能不悦。请注意，他是多么谨慎地藉着提及他们的方式，既荣耀了他们，又不伤害他。因为\Quote{教会}这个名称不会让主人发怒，即使他们与仆役被算在一起。因为教会不知道主仆之分；她以善行和罪恶来区分二者。如果那是一个教会，那么不要因你的奴隶与你一同受问候而不悦。因为\BibleQuote{在基督耶稣内，不再有奴隶或自由的人。}\BibleRef{迦 3:28}\par

\BibleQuote{愿恩宠与平安赐与你们。}\par

通过提及\Quote{恩宠}，他使他想起自己的罪过。他说：请考虑，天主在你身上赦免了多少大事，你是如何因恩宠得救的。效法你的主吧。他又为他祈求\Quote{平安}；这是理所当然的，因为当我们效法祂时，平安就会来临，恩宠也会常存。因为那对同仆不仁慈的仆人，直到他向同伴索要一百银钱时，仍蒙受主人的恩宠；但当他提出那个要求时，恩宠被夺走，他被交给刑役者。\par

道德教训。既然如此，让我们也向那得罪我们的人慈悲宽恕。我们这里的冒犯是一百银钱，但我们对天主的冒犯是一万塔冷通。但你知道，冒犯也因身份不同而受审判：例如，侮辱平民的人犯错了，但不如侮辱长官的人严重；冒犯更高长官的人有更大过错，冒犯较低长官的人过错较小；但冒犯君王的人罪大恶极。伤害本身相同，但因受辱者的卓越而变得更大。如果侮辱君王的人因君王的高位而受无法忍受的刑罚，那么侮辱天主的人要负多少塔冷通的债？因此，即使我们对天主犯下同样的罪，如同对人一样，那也不相等；因为天主与人的差别多大，对祂与对他们的冒犯之间的差别也多大。\par

但我如今也发现，过犯不仅因人物之卓越而变得重大，更是因其本身的性质。我将要说的话是可怕的，确实令人恐惧，但必须说出来，好让我们因之而震撼心灵、心生畏惧，表明我们敬畏人远胜于敬畏天主，尊崇人远胜于尊崇天主。试想，一个犯奸淫的人知道天主看见他，却不顾天主；但若有人看见他，他就抑制自己的情欲。这样的人岂不是不仅尊崇人超过天主，不仅侮辱了天主，而且——更可怕的是——他虽畏惧人，却藐视天主吗？因为他若看见人，就抑制情欲的火焰——但哪有什么火焰？那根本不是火焰，而是恣意妄为。倘若与女人相交本是不合法的，那也许可以称为火焰；但如今这却是侮辱和放纵。因为他若看见人，就会停止他的狂乱；而对天主的宽容，他却毫不在意。再者，一个行窃的人，明知自己在抢劫，却竭力欺骗人，为自己辩护，以漂亮言辞掩饰自己的托词；然而他无法如此欺瞒天主，却不在乎天主，不敬畏天主，不尊崇天主。如果国王命令我们不动用他人财物，甚至奉献自己的财物，大家都乐意捐献；但天主命令不可抢劫、不可夺取他人财物时，我们却不遵守。\par

你们看见了吗？我们尊崇人超过天主。这是个悲伤而令人痛苦的说法，是个沉重的控告。但请表明它是沉重的，从而逃离这事实！然而若你们不惧怕这事实，当你们说\Quote{我们害怕你的话，你给我们加上了重担}时，我怎能相信你们呢？是你们自己藉着行为给自己加上了重担，而不是我们的言语。若我仅仅提及你们所行之事的话语，你们就恼怒。这岂不荒谬吗？\par

但愿我所说的话是假的！我宁愿在那一天自己承担言语不当的指责，如同我徒然无因地责备了你们，也不愿见你们被控告这样的事。\par

你们不仅尊崇人超过天主，还强迫他人也如此行。许多人强迫他们的家仆和奴隶如此行。有些人违背他人的意愿强迫他们结婚，另一些人强迫他们服事可耻的差事、不洁的爱情、掠夺、欺诈和暴力的行为。因此这控告是双重的，他们也不能以必要性为由获得宽恕。因为即使你们自己不情愿地行恶事，出于统治者的命令，这也绝不是充分的借口；而当你们也强迫他人陷入同样的罪恶时，过犯就更重了。这样的人还有什么可宽恕的呢？\par

我说这些话，并非出于定你们罪的愿望，而是为表明我们在多少事上是天主的债户。因为我们若因尊崇人甚至与天主同等而侮辱了天主，那么当我们尊崇人超过祂时，又将如何？但如果那些对人所犯的过犯被证明是对天主更大的冒犯，那么当实际的过犯在其自身性质上更大更严重时，更是如此。\par

让每个人反省自己，他就会看见他所有的一切都是因为人。如果我们为天主的缘故所作的事，如同我们为人的缘故、为人的意见、为对人的恐惧或尊敬所作的事一样多，我们就是极其有福的了。如果我们有这么多要负责的事，我们应当极其乐意宽恕那些伤害我们、欺骗我们的人，并不存怨恨。因为有一条通向罪过赦免的路，不需要劳苦，不需要耗费财富，也不需要其他任何东西，只需我们自己的意愿。我们无需起程远行，无需出国界，无需冒危险和劳苦，只需愿意。\par

请问，在看似困难的事上，我们连一件轻省的事都不去做，而这又伴随着如此多的利益和益处，且毫无麻烦，我们还有什么借口？你不能轻看财富吗？你不能把你的财产分施给穷人吗？你不能愿意任何善事吗？你不能宽恕那得罪你的人吗？因为即便你没有这么多要负责的事，而天主只命令你宽恕，你难道不应该做吗？但如今你既然有这么多要负责的事，你还不宽恕吗？更何况你知道，你被要求这样做，是出于你从祂所领受的一切。如果我们去见我们的债主，他知情，就客气地接待我们，向我们表示尊敬，并大方地给予我们各种款待；虽然他并没有偿还债务，而是因为他希望我们在追债时变得仁慈。而你们这些欠天主如此多的人，被命令宽恕以期获得回报，难道你们不宽恕吗？请问，为什么不？我实在悲哀！我们领受了何等多的恩惠，却又以何等邪恶来回报！何等困倦！何等懈怠！德行何等容易，又伴随着如此多的益处；而邪恶何等劳苦！但我们却逃避那如此轻省的，去追求那比铅还重的。\par

这里不需要身体的强壮，不需要财富、产业、权势、友谊或其他任何东西；只要愿意，一切就成就了。有人使你忧伤、侮辱你、嘲笑你了吗？但请想想，你多少次对他人、甚至对主自己做了这样的事；你要容忍，宽恕他。请想想你所说的：\BibleQuote{宽免我们的罪债，如同我们也宽免得罪我们的人。}\BibleRef{玛 6:13}请想想，如果你不宽恕，你将不能有信心地说这话；但如果你宽恕，你就当作是债来要求，这不是出于事物的本性，而是出于那赐予者的仁慈。一个宽恕同仆的人，怎能与那获得对主所犯罪过的赦免相比呢？然而我们确实领受了如此大的仁慈，因为祂富于慈悲和怜悯。\par

为要表明即使没有这些事，也没有赦免，你因宽恕而得利，试想这样的人有多少朋友，他的赞美如何被四处走动的人传扬，他们说：\Quote{这是一个好人，他容易和解，他不知怀恨，他刚被打击就痊愈了。}当这样的人遭遇任何不幸时，谁不可怜他？当他冒犯时，谁不宽恕他？当他向他人求恩时，谁不给予他？谁不愿意成为如此善良灵魂的朋友和仆人？是的，我恳求你，让我们为祂做一切事，不仅是为我们的朋友，我们的亲属，甚至是为我们的家仆。因为祂说：\BibleQuote{你们要停止恐吓，因为知道你们的主也在天上。}\BibleRef{弗 6:9}\par

如果我们宽恕邻人的过犯，我们的过犯也将被宽恕；如果我们施舍，如果我们谦卑。因为这也除去罪过。因为如果那个税吏仅仅因为说：\BibleQuote{天主，可怜我这个罪人吧！}\BibleRef{路 18:13}，就得以成义回去，更何况我们，如果我们谦卑痛悔，也必能获得丰盛的慈爱。如果我们承认自己的罪并定自己的罪，我们便从大部分污秽中得洁净。因为有许多洁净的方法。因此，让我们以各种方式与魔鬼作战。我说的没有一件是难行的，没有一件是沉重的。宽恕那得罪你的人，怜悯穷人，谦卑你的灵魂，即使你是一个大罪人，也能藉着这些方法除去你的罪本身并擦去其污点，从而获得天国。愿天主使我们众人，在此世藉着告解从一切罪污中洁净自己后，在那里获得在基督耶稣我们的主内所预许的福份，等等。\par

\SourceRecord{Translated by Philip Schaff.；Nicene and Post-Nicene Fathers, First Series；Vol. 13.；Edited by Philip Schaff.；Buffalo, NY: Christian Literature Publishing Co.,；1889.；<http://www.newadvent.org/fathers/23091.htm>.}

% structure: p0001 source=h1 target=section reason=page-title

\section{\WorkTitle{费肋孟书第二篇讲道}}

\BibleRef{费 1:4-6}\par

\Quote{\Emph{我感谢我的天主，在祈祷中时常提起你，听说你的爱与信德，就是你对主耶稣和众圣人所所有的。愿你信德的交流，因认识我们内的一切善事，在基督耶稣内，得以发生功效。}}\par

他没有一开始就请求恩惠，而是先称赞那人，因他的善行赞扬他，并以不小的证据表明他的爱——他常在祈祷中提起他——并说许多人都因他得到安慰，且他在一切事上都顺从和依从；然后他才把请求放在最后，这样特别使他羞愧。\newline{}因为如果别人得到他们所求的，那么保禄更应得到；如果他在别人面前来求，他配得，何况他在别人之后来求，并且为别人而非自己求呢？然后，为了避免有人以为他写信仅为此事，或有人说：\Quote{如果不是为了敖尼息摩，你是不会写信的}，请看他如何还指出他这封书信的其他原因：首先是表明他的爱，其次也是希望为他预备住处。\par

\Quote{听说，}他说，\Quote{你的爱。}\par

这真是奇妙，比亲临目睹更大。因为显然那爱已因过度而显明，甚至传到了保禄。然而从罗马到夫黎基雅的距离并不小。他似乎曾在那里，因为提到了阿尔希波。因为哥罗森的居民是夫黎基雅人，他写信给他们时说：\Quote{当这封信在你们中间宣读后，也要使它在劳狄刻雅人的教会中宣读，而你们也要宣读从劳狄刻雅来的信。} \BibleRef{哥 4:16} 而这是一座夫黎基雅的城市。\par

他说：\Quote{愿你信德的交流，因认识我们内的一切善事，在基督耶稣内，得以发生功效。} 你看，他先给予，后接受；在求恩之前，先赐予他自己更大的一份。\Quote{愿你信德的交流，} 他说，\Quote{因认识你们内的一切善事，在基督耶稣内，得以发生功效；} 也就是说，愿你达到一切德行，一无所缺。因为信德这样才发生功效，即有行为相随。因为\Quote{没有行为的信德是死的。} \BibleRef{雅 2:26} 而他没说\Quote{你的信德}， 而是\Quote{你信德的交流}， 将自己与他连接起来，表明是一个身体，并以此特别使他羞于拒绝。他说，如果你在信德上有分，你也应当在其他事上相通。\par

第7节。\Quote{因为我们因你的爱得了极大的喜乐和安慰，因为众圣人的心肠因你而得了舒畅，弟兄。} \BibleRef{费 1:7}\par

没有什么比提出向他人所行的善事更能使我们羞愧而施予，尤其当那人比他人更应受尊重时。他没有说\Quote{如果你行在别人身上，更当行在我身上}； 但他暗示了同一件事，尽管他以另一种更温和的方式做到了。\par

\Quote{我得了喜乐，} 这就是说，你从你向别人所行的事中给了我信心。\Quote{和安慰，} 这就是说，我们不仅感到愉快，而且得了安慰。因为他们是我们的肢体。既然应有如此的一致，以致当任何其他受苦的人得了舒畅，即使我们自己一无所获，也应当为他们高兴，如同一个身体受益；那么，你若也使我们得了舒畅，就更当如此。他没有说\Quote{因为你顺从和依从}， 而是更强烈、更强调地说\Quote{因为众圣人的心肠}， 如同父母对自己所疼爱的孩子，以致这爱和情感表明他也被他们极其爱慕。\par

第8节。\Quote{因此，我虽然在基督里能放胆吩咐你行那相宜的事。} \BibleRef{费 1:8}\par

请看他是何等谨慎，免得那些出于极爱而说的话，会如此刺激听众而使他受伤。因此，在他说\Quote{吩咐你}之前， 因这话是冒犯的——尽管出于爱，更应抚慰他——但出于极度的审慎，他好像纠正了这话，说\Quote{有了信心}， 由此暗示费肋孟是一位伟人，即\Quote{你给了我们信心。} 不仅如此，还加上\Quote{在基督里}这个说法，表明不是他在世上更显赫，不是他更有权力，而是因他对基督的信德——然后他又加上\Quote{吩咐你}， 不仅这样，还有\Quote{那相宜的事}， 即合理的行动。请看他从多少事物中为此提出证据：你对别人行善，也对我行善，并且是为了基督，又因那事是合理的，以及爱所给予的，所以他又加上：\par

第9节。\Quote{然而为爱的缘故，我宁愿恳求你。} \BibleRef{费 1:9}\par

好像他曾说：我知道我确实能因已发生的事，以极大的权柄命令而成就；但因为我对此事极为关心，\Quote{我恳求你。} 他同时表明这两点：他对费肋孟有信心，因为他命令他；又极其关心此事，所以他恳求他。\par

\Quote{像我这样的一位，} 他说，\Quote{即年老的保禄。} 奇怪！这里有多少事使他羞愧而顺从！保禄，从他的品性，从他的年龄——因他年老——和从比一切都更正义的事——他还成了\Quote{耶稣基督的囚犯}。\par

因为谁会不用张开双臂欢迎一位已经获冠的战士呢？谁看见他为基督的缘故被捆绑，不肯向他施予千万恩惠呢？他在此之前用如此多的考虑抚慰了他的心，没有立刻提出名字，而是推迟作出如此重大的请求。因为你们知道主人对逃跑的奴隶怀有何等心态；尤其当他们还带了盗窃时，即使遇到好主人，他们的怒气也会增加。他费尽苦功安抚了这怒气，先说服他乐意在任何事上服事他，并预备了他的灵魂去顺服一切，然后才提出请求，并加上了赞美，\Quote{我请求你，} \Quote{为我在锁链中所生的儿子。}\par

再次提到锁链是为使他羞愧而顺从，然后才提到名字。因为他不仅熄灭了怒气，还使他喜悦。因为他说：若他不是特别有益，我本不会称他为我的儿子。我所称呼弟茂德的那名字，我也这样称呼他。并且一再显示他的情感，用他重生的时期来催促他，\Quote{我在锁链中生了他，}他说，好叫因此他也值得获得极大的光荣，因为他是在他的战斗、他为基督所承受的考验中重生的。\par

\Quote{敖尼息摩，}\par

十一节。\Quote{他从前对你没有益处。} \BibleRef{费 1:11}\par

看他的智慧多么伟大，他如何承认那人的过错，从而熄灭了怒气。他说：我知道，他从前是无益的。\par

\Quote{但现在}他将\Quote{对你和我都有益处。}\par

他没有说他对你有益，以免引起反驳，而是引入了自己的身份，使他的希望显得可信，\Quote{但现在，}他说，\Quote{对你和我都有益处。}因为若他对如此要求严格的保禄有益，那么他对他的主人更有益。\par

十二节。\Quote{我现在打发他回你那里去。} \BibleRef{费 1:12}\par

由此他也熄灭了怒气，通过交出他。因为当主人被请求宽恕不在场的奴隶时，他们最是恼怒；所以通过这个举动，他更安抚了他。\par

十二节。\Quote{你务必收下他，他是我自己的心肠。} \BibleRef{费 1:12}\par

他又没有仅仅说出名字，而是附带了一个能感动他的话，这比儿子更亲。他说了\Quote{儿子}，说了\Quote{我生了他}，所以很可能他会非常爱他，因为他在考验中生了他。因为显然我们对于那些在自己脱离危险之后所生的孩子，感情最为炽烈，就如经上说的：\Quote{祸哉，巴罗加贝尔！}又如辣黑耳给本雅明起名\Quote{我苦难的儿子。} \BibleRef{创 35:18}\par

\Quote{所以你，}他说，\Quote{务必收下他，他是我自己的心肠。}他显示了他情感的深厚。他没有说：领回去，他没有说：不要生气，而是\Quote{收下他}；意思是：他不仅值得宽恕，还值得尊敬。为什么？因为他成了保禄的儿子。\par

十三节。\Quote{我本来有意将他留在我这里，叫他在你方面替我服事我在福音的锁链中。} \BibleRef{费 1:13}\par

你们看见吗？经过这么多的预备，他终于光荣地将他带到主人面前，并且看他是多么有智慧地做了这事。看他使他对多少事负责，又是多么尊重另一方。他说：你找到了一个借着他归还我服事的方法。在这里他显示了他更顾念主人的益处而不是奴隶的益处，并且非常尊敬他。\par

十四节。\Quote{但未得你的同意，}他说，\Quote{我什么也不愿做；好叫你的善行不是出于勉强，而是出于甘心。} \BibleRef{费 1:14}\par

这一点特别使被请求的人感到荣幸：当事情本身有益，并且是在他同意下提出时。因为这产生两个好效果：一方获益，另一方更稳固。他没有说：免得出于勉强，而是\Quote{像是出于勉强}。因为我知道，他说，你并非未经知晓，而是一下子就得知了，你不会生气，但出于过分的考虑，使它\Quote{不像是出于勉强}。\par

十五至十六节。\Quote{因为他暂时离开你，也许正是为此，好叫你可以永远得到他；不再是奴仆。} \BibleRef{费 1:15-16}\par

他巧妙地说\Quote{也许}，好让主人让步。因为逃跑出于邪恶和败坏的意念，并非出于这种意图，所以他说\Quote{也许}。他没有说：所以他逃跑了，而是：所以他\Quote{被分开}，用更好听的词进一步安抚他。他没有说：他分开了自己，而是\Quote{他被分开}。因为他离开并非出于自己的安排，既不为这个目的也不为那个。这也是若瑟为他弟兄们辩护时说的：\Quote{天主差我到这里来} \BibleRef{创 45:5}，也就是说，天主利用了他们的恶达到善的目的。\Quote{所以，}他说，\Quote{他暂时被分开。}这样他缩短了时间，承认了过犯，并将一切都转向天主的眷顾。\Quote{好叫你可以得到他，}他说，\Quote{永远，}不只是现在，甚至将来，好叫你永远拥有他，不再是个奴隶，而是比奴隶更尊贵。因为以后他不会再逃跑。\Quote{好叫你可以得到他，}他说，\Quote{永远，}也就是说，重新拥有他。\par

\Quote{不再是奴仆，而是高于奴仆，是亲爱的弟兄，尤其对我。} \BibleRef{费 1:16}\par

你失去一个奴隶不过是短暂的，但你将永远得到一个弟兄，不仅是你自己的弟兄，也是我的弟兄。在这里也有很大的德行。但如果他是我的弟兄，你也不会以他为耻。称他为自己的儿子，显示了他的骨肉之情；称他为弟兄，显示了他对他的极大善意，以及与他同等的尊荣。\par

道德：这些事并非没有目的而记下来，而是为了使我们作主人的不要绝望于我们的仆人，也不要过分压迫他们，而要学习宽恕这类仆人的过失，使我们不总是严厉，也不因他们的奴仆身份而羞于在他们善良时，使他们与我们共享一切。因为如果\Person{保禄}不以称一人为\BibleQuote{他的儿子，他自己的心肠，他的弟兄，他所爱的}为耻，那么我们当然也不该以此为耻。我为什么说\Person{保禄}呢？\Person{保禄}的主不以称我们的仆人为他自己的弟兄为耻，而我们却羞耻吗？看祂如何尊崇我们；祂称我们的仆人为自己的弟兄、朋友和共同承嗣者。看祂降卑到了什么地步！那么，我们做了什么，就完成了我们的全部责任呢？我们永远无法做到；但无论我们达到何种谦逊的程度，还有更大的部分留在后面。因为试想，无论你做什么，你都是在对待一个同仆，而你的主却向你的仆人做了这些。听啊，战栗吧！绝不要因你的谦逊而自高自大！\par

也许你会嘲笑这种说法，好像谦逊会使人自高自大。但不要为此惊讶，当它不是真正的谦逊时，它确实会使人自高自大。怎么一回事呢？当它只是为了赢得人的欢心，而不是天主的欢心，为了被人称赞，变得高傲时。因为这也就是魔鬼的行径。正如许多人为不虚荣而虚荣，同样他们也因谦卑自己而骄傲，因为他们是高傲的。例如，一个弟兄来了，甚至一个仆人来了，你接待了他，洗了他的脚；你立刻自视甚高。你说，我做了别人没有做的事。我达到了谦逊。那么，一个人如何才能保持谦逊呢？如果他记住基督的命令，祂说：\BibleQuote{你们做了一切事，要說：我们是无用的仆人。}\BibleRef{路 17:10}以及世间的导师又说：\BibleQuote{我却不认为我已经得着了。}\BibleRef{斐 3:13}凡说服自己没有做任何大事的人，无论他做了多少事，只有他能谦卑自下，认为自己尚未达到完美的人。\par

许多人因他们的谦逊而骄傲；但我们不要受此影响。你做了什么谦逊的事吗？不要为此自豪，否则它的全部功绩都失去了。法利塞人就是这样，他因将自己所得的十分之一施舍给穷人而自高自大，并失去了所有的功绩。\BibleRef{路 18:12}但税吏却不然。再听\Person{保禄}说：\BibleQuote{我虽自觉良心无愧，但我却并不因此而成义。}\BibleRef{格前 4:4}你看见他没有高举自己，而是千方百计地贬抑自己，谦卑自己，而且是在他已经登上了顶峰的时候。三位青年在火里，在火窑中，他们说了什么？\BibleQuote{我们和我们的祖先犯了罪，行了不义。}\BibleRef{歌 5:6}这就是拥有痛悔之心；为此他们可以说：\Quote{但愿我们在痛悔的心和谦逊的精神里蒙悦纳。}因此即使他们掉进火窑之后，他们也极其谦卑，甚至比以前更加谦卑。因为当他们看到所行的奇迹时，自认为不配得那拯救，就更加谦卑。因为当我们确信自己蒙受了不该得的大恩时，我们就特别悲伤。然而他们得到了什么不应得的恩惠呢？他们把自己交给了火窑；他们因别人的罪而被掳掠；因为他们还年轻；他们没有抱怨，也没有愤慨，也没有说：我们事奉天主有什么好处？敬拜祂有什么益处？这个人不虔敬，却成了我们的主。我们与拜偶像的人一同被一个拜偶像的君王惩罚。我们被掳掠了。我们失去了祖国、自由、一切祖产，我们成了囚犯和奴隶，我们被奴役于一个野蛮的君王。他们没有说这些话中的任何一句。但说了什么？\Quote{我们犯了罪，行了不义。}他们不仅为自己，也为别人献上祈祷。因为他们说：\Quote{你将我们交在一个可憎的、邪恶的君王手里。}再者，\Person{达尼尔}第二次被投入狮子坑时，说：\Quote{因为天主记念了我。}为什么祂不该记念你呢，\Person{达尼尔}？当你在君王面前光荣祂说：\BibleQuote{不是因为我有智慧}\BibleRef{达 2:30}？但当你因为没有服从那最邪恶的法令而被投入狮子坑时，为什么祂不该记念你呢？当然正是为此祂应该记念你。你不是为了祂的缘故而被投入的吗？\Quote{是的，}他说，\Quote{但我欠了许多债。}如果他在显示了如此巨大的德行之后还说这样的话，那么在这之后我们该说什么呢？但听\Person{达味}说：\BibleQuote{他若这样说：我不喜悦你，看，我在这里，愿他照看为好的待我。}\BibleRef{撒下 15:26}然而他有无数好事可说。\Person{厄里}也说：\BibleQuote{这是上主，愿他行他看为好的事。}\BibleRef{撒上 3:18}\par

这是善仆的本分，不仅在天主的仁慈中，也在祂的责罚和惩罚中完全顺服于祂。如果我们忍受主人责打仆人，知道他们会饶恕他们，因为他们是自己的仆人，却以为天主在惩罚时不会饶恕，这岂不是荒谬的吗？\Person{保禄}也暗示了这一点，说：\BibleQuote{我们无论活，无论死，都是属于主的。}\BibleRef{罗 14:8}我们说，一个人不希望自己的财产减少，他知道如何惩罚，他在惩罚自己的仆人。但确实，没有比祂更宽恕我们的了，祂从无中创造了我们，祂使太阳升起，祂降雨，祂将生命吹入我们内，祂为我们赐下自己的儿子。\par

但正如我先前所说，也是我之所以说这一切的原因，让我们按应当的怀有谦卑之心，按应当的节制。不要让它成为我们自高自大的借口。你谦卑，且比众人更谦卑吗？不要因此自视过高，也不可指责他人，免得你失去夸口。因为你谦卑正是为了从骄傲的狂乱中得释放；因此，若通过你的谦卑你陷入那种狂乱，那倒不如你不谦卑。请听保禄说：\Quote{罪恶借着那善的在我内产生死亡，为使罪恶借着诫命成为极其罪恶的。}\BibleRef{罗 7:13}当你想因自己谦卑而崇拜自己时，思想你的主降生成人至何等境地，你便不再崇拜自己，也不赞美自己，反会嘲笑自己一无所成。看自己完全是个负债者。无论你做了什么，要记得那比喻：\Quote{你们中间谁有仆人……他进来时，会对他说：你来坐下吃饭吗？……我告诉你们：不……但你要侍候我。}\BibleRef{路 17:7}我们岂向侍候我们的仆人致谢吗？绝不。然而天主却感谢我们，我们并非侍候祂，而是做对自己有益的事。\par

但我们不要这样想，仿佛祂欠我们感谢，好使祂更欠我们，而要仿佛我们是在偿还债务。因为这事确实是债务，我们所做的一切都是债务。因为若我们用金钱购买奴隶，就希望他们完全为我们而活，凡他们所有的都归我们所有，更何况是那位从无中创造我们、又以祂的宝血赎买了我们、为我们付上了无人愿为自己的儿子付上的代价、为我们倾流祂自己宝血的主呢？因此，即使我们有千万条灵魂，并为祂全都舍弃，我们岂能对祂做出相等的回报？绝不。为什么？因为祂这样做，并不欠我们什么，全是恩宠。但我们从此是负债者：祂本是天主，却成为仆人；祂本不受死亡支配，却在肉身上甘受死亡。我们若不为其献出生命，按自然律也必定要交出，稍后也必不情愿地与生命分离。在财富上也是如此：若不为其奉献，临终时也必不得已交出。在谦卑上亦然：即使不为其谦卑，也必因磨难、灾祸、掌权者而变得谦卑。你看，这是何等大的恩宠！祂没有说：\Quote{殉道者做了何等大事？他们即使不为我而死，也必定会死。}但祂却认为自己大大欠了他们，因为他们自愿放弃按自然律不久将不情愿放弃的。祂没有说：\Quote{施舍财富者做了何等大事？他们即使不情愿，也必得交出。}但祂却认为自己大为亏欠他们，并且不羞于在众人面前承认，主竟由祂的仆役供养。\par

因为这也是一位主人的荣耀：拥有感恩的仆役。这也是一位主人的荣耀：如此爱祂的仆役。这也是一位主人的荣耀：将属于仆役的视为己有。这也是一位主人的荣耀：不羞于在众人面前承认他们。因此，让我们因基督如此伟大的爱而心生敬畏。让我们被这爱之剂点燃。即使一个人卑微低贱，但我们若听说他爱我们，就必在一切事上因对他的爱而火热，并极其尊重他。我们尚且如此爱，而我们的主如此爱我们，我们岂不激动？我恳求你们，不要对我们灵魂的救恩如此冷漠，而要按我们的能力爱祂，并将一切都用于祂的爱：我们的生命、财富、光荣，一切都要怀着喜乐、欢喜、热忱，不是作为向祂的回报，而是为我们自己。因为这是爱者的法则：他们以为在为自己所爱者受苦时，是领受恩惠。因此，让我们以这样的心情对待我们的主，好使我们也能分享在基督耶稣我们的主内将来的美善。\par

\SourceRecord{Translated by Philip Schaff.；Nicene and Post-Nicene Fathers, First Series；Vol. 13.；Edited by Philip Schaff.；Buffalo, NY: Christian Literature Publishing Co.,；1889.；<http://www.newadvent.org/fathers/23092.htm>.}

% structure: p0001 source=h1 target=section reason=page-title

\section{《费肋孟书》第三篇讲道}

费肋孟书 一章17至19节 \BibleRef{费 1:17-19}\par

\Quote{你若以我为伙伴，就收留他如同收留我一样。他若亏负了你，或欠你什么，都记在我的账上。我保禄亲笔写下：我必偿还。我并不对你说，连你自己也亏欠于我。}\par

没有比不一下子要求一切更易于使人听从的了。请看，经过多少赞美，多少准备之后，他才提出这件大事。说了他是\Quote{我的儿子}、是福音的分享者、是\Quote{我的心肠}、你收留他\Quote{如同兄弟}、并\Quote{以他为兄弟}之后，他又加上\Quote{如同我一样}。保禄并不以此为耻。因为那不以被称为信友的仆人为耻、而且承认自己就是这样的人，更不会拒绝这事。他所说的话大意是：你若与我同心，若在同样的条件下奔走，若以我为朋友，就收留他如同我一样。\par

\Quote{他若亏负了你。}请看他在何时何处提到伤害的事：是在说了许多为他辩护的话之后。因为金钱的损失特别容易使人烦恼，为了避免他为此责备他（很可能钱已花掉），便提出这事说：\Quote{他若亏负了你。}他不是说：他若偷了什么；而是：\Quote{他若亏负了你。}他同时既认了错，又不把它当作仆人的错，而是朋友对朋友的错，用了\Quote{亏负}一词而非\Quote{偷窃}。\par

\Quote{都记在我的账上，}他说，意思是，算作欠我的债，\Quote{我必偿还。}然后带着那种属灵的诙谐语气，\par

\Quote{我保禄亲笔写下。}既感人又诙谐：如果保禄不拒绝为他立下字据，他倒拒绝收留他！这既使费肋孟惭愧而顺从，又使敖尼息摩脱离困境。他说：\Quote{我写下，}\Quote{以我自己的手。}没有什么比这些\Quote{心肠}更充满深情，没有什么比他更热切，更热诚的了。请看，他为一个人付出了何等的关怀。然后，为了不让被请求的人觉得难堪，好像他没有胆量为一件盗窃之事求情并成功，他稍微缓和一下，说：\Quote{我并不对你说，连你自己也亏欠于我。}不单是你的东西，连你自身也亏欠我。这是出于爱，符合友谊的规则，也是他极大信任的证明。请看，他如何处处兼顾，既大胆请求，又不显得对对方信任不足。\par

20节 \BibleRef{费 1:20} \Quote{是的，兄弟。}\par

\Quote{是的，兄弟}是什么意思？收留他，他说。这一点我们必须从上下文理解，虽未明说。因为他放下诙谐，又回到先前的严肃话题。然而这些也是严肃的，因为圣者所发之言，即使是诙谐，本身也是严肃的。\par

\Quote{是的，兄弟，让我从你身上因主而喜乐，使我的心在基督内得安息。}\par

也就是说，你是为主行这善事，并非为我。\Quote{我的心}指我对你的心。\par

21节 \BibleRef{费 1:21} \Quote{我深信你的服从，才写信给你。}\par

还有哪块石头不被这些话软化？还有哪头野兽不被这些话驯服，并准备全心接待他？在用了那么多巨大的见证称赞他的良善之后，他并不耻于再次为自己辩解。他说：并非仅仅请求，也不是命令，更不是武断，而是\Quote{深信你的服从，才写信给你}。他在信的开头说了\Quote{深信}，在信的结尾也这样说。\par

\Quote{知道你会做得比我说得更多。}\par

同时，他说这话也激励了他。因为若他在这方面有如此信誉，即他会做得比说得更多，而他若不做到这么多，他自己也会惭愧。\par

22节 \BibleRef{费 1:22} \Quote{同时，也请为我预备住处，因为我希望借着你们的祈祷，我会被赐给你们。}\par

这也是极其信任的表现——不如说，这也是为了敖尼息摩，好叫他们不漠不关心，而知道他在返回后会知道有关他的事，从而放下对过错的记忆，更乐意施恩。因为保禄住在他们中间，保禄老年时，保禄出狱后，他的影响力和尊荣是巨大的。再说，他说他们祈祷，这是他们爱德的证明，并归功于他们，好像他们为他祈祷一样。因为，他说，虽然我现在身处危险，但你们若为此祈祷，仍会见到我。\par

23节 \BibleRef{费 1:23} \Quote{厄帕夫辣，我在基督耶稣内的同囚，问候你。}\par

他是哥罗森人派遣来的，由此可见费肋孟也在哥罗森。他称他为\Quote{同囚}，表明他也经受了许多磨难，因此即使不为他自己，也该为他人的缘故而应允所求。因为那在磨难中不顾自己而关心他人的人，是值得应允的。\par

他又从另一方面使他惭愧：如果他的同乡与保禄一同被囚，与他一同受苦，而他竟不为自己的仆人向保禄施恩。他又加上\Quote{我在基督耶稣内的同囚}，意思是\Quote{为基督的缘故}。\par

24节 \BibleRef{费 1:24} \Quote{马尔谷、阿黎斯塔苛、德玛斯、路加，我的同工们，问候你。}\par

那么为何他将路加列在最后？但他别处说：\Quote{\BibleQuote{只有路加与我同在}}\BibleRef{弟后 4:11}，又说：\Quote{\BibleQuote{德玛斯}}是那些\Quote{\BibleQuote{离弃了他，因贪爱现世}}的人之一。\BibleRef{弟后 4:10}所有这些事，虽在别处已提及，但仍然不可不经查究便在此略过，也不可只视作寻常之事而听之。但他怎会说离弃他的人向他们致候呢？因为他说：\Quote{\BibleQuote{厄辣斯托留于格林多。}}\BibleRef{弟后 4:20}他又加上厄帕夫辣，因他们认识他，且是他的同乡。马尔谷亦同，因他自己也是一位可敬之人。那么他为何将德玛斯与这些人并列？也许此后他见到危险倍增，便更懈怠了。但路加虽居末位，却成了首位。他以此等人向他致候，更激励他服从，并称他们为同工，以此羞愧地迫使他应允所求。\par

第25节。\Quote{\BibleQuote{愿我们的主耶稣基督的恩宠与你的心灵同在。阿们。}}\BibleRef{费 1:25}\par

伦理教训。他以祈祷结束书信。祈祷确实是大善，有益于灵魂的得救与保存。但当我们行配得之事，不使自己成为不配时，它才是伟大的。因此，当你去见司铎，他若对你说：\Quote{主必怜悯你，我儿}，你不可只依赖话语，也要加上行动。行堪受怜悯之事，天主必祝福你，我儿，只要你行配得祝福之事。你若向近人施怜悯，祂必祝福你。因为我们愿从天主获得的，应先施予近人。若我们剥夺近人这些，又怎能期望获得呢？祂说：\Quote{\BibleQuote{怜悯人的人是有福的，因为他们要受怜悯。}}\BibleRef{玛 5:7}若人对这样的人施怜悯，何况天主呢？但对不怜悯者，决不施怜悯。\Quote{\BibleQuote{因为那不怜悯人的，也要受无怜悯的审判。}}\BibleRef{雅 2:13}\par

怜悯何其美！你为何不对他人行呢？你得罪时希望被宽恕，为何自己不宽恕得罪你的人？你来天主前，求祂赐天国，而你自己在别人乞求时却不肯施舍银钱。因此我们不获怜悯，因为我们不施怜悯。为何如此？你说。对不怜悯的人施怜悯，不也是怜悯的一部分吗？不然！对那待邻人心硬冷酷、作恶无数之人，若施以极大仁慈，怎能算是怜悯呢？那么，你说，那洗净我们无数罪恶的圣洗岂不拯救了我们？它释放了我们，不是要我们再犯，而是使我们不再犯。因为经上说：\Quote{\BibleQuote{我们已死于罪恶，怎能还生活在罪中呢？}}\BibleRef{罗 6:2}\par

\Quote{\BibleQuote{那么，我们因不在法律之下，就可以犯罪吗？绝对不可！}}\BibleRef{罗 6:15}为此，天主解救你脱离那些罪恶，使你不再退回到那耻辱中。就如医生解除病人的热症，不是要他们滥用健康，自招损伤混乱（因为若一人用健康只为使自己再度卧床，则生病倒更好），而是为使他们在明白疾病之害后，不再陷入同样状况，更稳妥地保全健康，并力行一切有益健康之事。\par

那么，你说，若天主不拯救恶人，祂的慈爱何在？因为我常听许多人这样说：祂是人的朋友，必会拯救众人。为避免我们徒然自欺（因我记得曾向你们许下这类诺言），来，今天让我们就此论点展开讨论。我最近与你们讲论地狱，并延缓了关于天主慈爱的论证。因此今天适宜重拾此话题。地狱将存在，我们已充分证明，举出洪水及从前的灾祸，并论证那行这些事的天主，不可能不惩罚现今世代的人。若祂惩罚了律法以前犯罪的人，祂也不会放过那些在恩宠之后犯更大罪恶的人不受罚。于是有人质疑：祂怎会是善？怎会怜悯人，如果祂至少惩罚人的话？我们延缓了论证，以免用过多言语淹没你们的耳朵。\par

来，今天让我们还清债务，并表明天主即使在惩罚中也是善的。因为这篇言论适合我们反对异端者。让我们认真注意。天主不需求我们任何事物，却造了我们。祂不需求我们，从祂在长时间后才造我们便可证明。若祂需要我们，祂早就造了我们。因为祂自己存在，即使没有我们，而我们在后世才被造，祂造我们并非因为需要。\par

祂造了天、地、海及一切存在之物，皆为我们。告诉我，这难道不是善的标记吗？可有许多事可提。简言之，\Quote{\BibleQuote{祂使太阳上升，光照恶人，也光照善人；降雨给义人，也给不义的人。}}\BibleRef{玛 5:45}这难道不是善的标记吗？不，你答。因我曾与一马西翁派信徒交谈时问：这些岂非善的标记？他答：若祂不追究人的罪，那才是善。但若祂追究罪，就不是善。此人今不在场。但来，让我们重述当时所讲，并再加些话。我因丰裕而表明：若祂不追究人的罪，祂就不是善；正因祂追究罪，所以祂是善。\par

比如说，如果祂不向我们追究，那时人的生命还能存续吗？\newline{}那时我们岂不陷入了野兽的状态吗？\newline{}因为如果当这恐惧、这追究和审判悬在我们头上时，我们已在彼此吞吃上超过了鱼类，在掠夺彼此的财产上使豺狼和狮子相形见绌；\newline{}如果祂不向我们追究，而我们也确信这一点，那么生命将被何等大的骚动和混乱充满呢？\newline{}此后，传说中的迷宫与世界上的纷乱相比，又算得了什么呢？\newline{}你岂不是会看到无数的淫秽和混乱？\newline{}因为那时谁还会再尊敬他的父亲？谁会饶恕他的母亲？谁会不尝试任何快乐、任何邪恶呢？\newline{}事情就是这样，我要设法只从一个家庭向你展示。如何呢？\newline{}你们这些提出这些问题并拥有仆人的人；如果我能向这些人显明，如果他们毁坏主人的家庭，侮辱主人的身体，掠夺一切，把一切搞得天翻地覆，对待主人如同敌人，而主人却不威胁他们、不纠正他们、不惩罚他们，甚至不用一句话使他们难过，这难道是良善的证据吗？\newline{}我认为这是极端的残忍，不仅因为这不合情理的仁慈出卖了妻子和儿女，也因为这些奴隶自身在他们面前被毁灭。\newline{}因为他们将成为酒徒、放荡、荒淫，比任何野兽都更无理性。\newline{}告诉我，这难道是良善的证据吗？践踏灵魂的高贵本性，毁灭他们自己和别人？\newline{}你难道看不出，追究人的责任是伟大良善的证据吗？\newline{}但我为什么提到奴隶，他们更容易陷入这些罪恶？\newline{}但让一个人有儿子，任凭他们为所欲为，而不惩罚他们；他们岂不比任何事物更坏吗？\newline{}告诉我。那么在人中间，惩罚是良善的标志，不惩罚是残忍的标志，在天主身上难道不是这样吗？\newline{}所以，正因为祂是善的，因此祂预备了地狱。\par

你还希望我讲另一个天主良善的例子吗？这不仅在于此，还在于祂不让善人变为恶人。因为如果他们注定遭遇同样的事情，他们都会变坏。但现在这一点也给善人不少安慰。因为请听先知所说：\Quote{义人见对不敬虔者复仇时，必欢喜，他要在罪人的血中洗手。} \BibleRef{咏 57:10}\newline{}不是因复仇而欢喜，愿天主禁止！而是害怕自己遭受同样的事，他会使自己的生命更加纯净。这便是祂极大关怀的标志。\newline{}是的，你说，但祂本应只警告，而不应也惩罚。但如果祂真的惩罚，而你仍说这只是警告，并因此变得更懒散，如果它真的只是警告，难道你不会变得更懈怠吗？\newline{}如果尼尼微人知道那只是警告，他们就不会悔改。但因为她们悔改了，她们使警告只停留在言语上。你希望它只是警告吗？你对此有处理之权。变成一个更好的人，它就只停在警告。但如果你——愿此事远离你！——轻视这警告，你将会经历它。\newline{}洪水之前的人们，如果曾畏惧那警告，就不会经历它的执行。而我们，若畏惧那警告，就不至于使自己经历事实。愿天主禁止我们如此。愿仁慈的天主恩赐我们所有人从此以后，得获健全心智，获得那些不可言喻的福分。愿我们都堪当承受这些福分，藉着我们的主耶稣基督的恩宠和慈爱，愿光荣、权能和尊崇归于祂及圣父和圣神，从现在直到永远。阿们。\par

\SourceRecord{Translated by Philip Schaff.；Nicene and Post-Nicene Fathers, First Series；Vol. 13.；Edited by Philip Schaff.；Buffalo, NY: Christian Literature Publishing Co.,；1889.；<http://www.newadvent.org/fathers/23093.htm>.}
